\subsection{Sparse replacements on prescribed vertex sets}
\label{app:factor-expanders}

We use the Margulis--Gabber--Galil family from
Theorem~\ref{thm:auxiliary-expander-family} and the expander mixing
lemma (Lemma~\ref{lem:expander-mixing}).

Fix $\zeta>0$ and $\kappa\ge1$ as in Lemma~\ref{lem:factor-expanders}.
We map the vertices of a powered expander onto the prescribed sets
and verify the degree, no-hole, and connectedness guarantees.

Use $\lambda_0=5\sqrt2/8$ from
Theorem~\ref{thm:auxiliary-expander-family}. Fix an integer $h\ge1$
with $\lambda_0^h<\zeta/4$, and set
$d_{\zeta,\kappa}=\lceil4\kappa\rceil8^h$.
Use this same walk length for both constructions. For each integer
$\nu\ge1$, put $M=\nu^2$ and take
$\mathcal E=\mathcal G_\nu^{[h]}$ from that theorem.
It has degree $d=8^h$ and normalized operator norm at most
$\lambda_0^h$ on the zero-sum subspace, so its nonconstant
adjacency eigenvalue bound is $d\lambda_0^h$.
For $S,T\subseteq V(\mathcal E)$, Lemma~\ref{lem:expander-mixing} gives
\begin{equation}\label{eq:factor-expander-mixing}
 e_{\mathcal E}(S,T)
 \overset{\text{Lem.~\ref{lem:expander-mixing}}}{\geq}
 d\sqrt{|S||T|}\left(\frac{\sqrt{|S||T|}}{M}-\lambda_0^h\right).
\end{equation}
For an integer $t\ge1$, if $|S|,|T|\ge t>\lambda_0^hM$, then
$\sqrt{|S||T|}/M\ge t/M>\lambda_0^h$, so the right-hand side is
positive. Hence a step joins $S$ to $T$.

We accommodate arbitrary sizes by grouping auxiliary vertices and
assigning each group to a prescribed vertex. Auxiliary steps give
edges between distinct images; group sizes bound the output degrees.

For a map $\pi:V(\mathcal E)\to U$ and a subset $X\subseteq U$,
the \emph{preimage} is
$\pi^{-1}(X)=\{x\in V(\mathcal E):\pi(x)\in X\}$.
The map is \emph{onto} if every vertex of $U$ has at least one
preimage. When $1\le n=|U|\le M$, partition the $M$ auxiliary vertices
into $n$ groups. Write $M=n\lfloor M/n\rfloor+r$, where
$r$ is an integer with $0\leq r<n$. Give $r$ groups size $\lfloor M/n\rfloor+1$ and the
other groups size $\lfloor M/n\rfloor$. Since $M\geq n$, every group
is nonempty. Assign one group to each vertex of $U$. This gives
\begin{equation}\label{eq:factor-expander-preimages}
 \lfloor M/n\rfloor\le|\pi^{-1}(\{u\})|\le\lceil M/n\rceil
 \qquad(u\in U).
\end{equation}

Use the following two constructions.
\begin{enumerate}
 \item\label{constr:one-set-replacement} \emph{One prescribed set.}
 Let $U$ have size $n\ge1$, put $M=\lceil\sqrt n\rceil^2$, and
 take the $M$-vertex graph $\mathcal E$ above. Choose an onto map
 $\pi:V(\mathcal E)\to U$ with the group sizes just specified.
 Define the replacement $R_U$ by
 \[
 V(R_U)=U,
 \qquad E(R_U)=\{uv\in\tbinom{U}{2}:
 e_{\mathcal E}(\pi^{-1}(\{u\}),\pi^{-1}(\{v\}))>0\}.
 \]
 Thus equal images give no edge, and repeated pairs give a single edge.
 \item\label{constr:two-side-replacement} \emph{Two prescribed sides.}
 Let $A,B$ be disjoint sets of sizes $n_A,n_B\ge1$, put
 $M=\lceil\sqrt{\max\{n_A,n_B\}}\rceil^2$, and take the
 $M$-vertex graph $\mathcal E$ above. Choose onto maps
 $\pi_A:V(\mathcal E)\to A$ and
 $\pi_B:V(\mathcal E)\to B$, with preimage sizes
 $\lfloor M/n_A\rfloor$ or $\lceil M/n_A\rceil$ for the first
 map, and $\lfloor M/n_B\rfloor$ or $\lceil M/n_B\rceil$ for
 the second. Define the replacement $R_{A,B}$ by
 \[
 V(R_{A,B})=A\mathbin{\dot\cup}B,
 \qquad E(R_{A,B})=\{uv:u\in A,\ v\in B,
 e_{\mathcal E}(\pi_A^{-1}(\{u\}),\pi_B^{-1}(\{v\}))>0\}.
 \]
 This gives one edge for each distinct pair and no edge within
 either side. The two maps have the same
 domain, and their preimage groups may differ.
\end{enumerate}

A walk is \emph{closed} if its first and last vertices coincide.

\begin{restatedlemma}{Lemma~\ref{lem:factor-expanders}}
For every fixed $\zeta>0$ and $\kappa\geq1$, there is a constant
$d_{\zeta,\kappa}$ such that the following connected simple graphs,
each of maximum degree at most $d_{\zeta,\kappa}$, can be constructed
deterministically in polynomial time.
\begin{enumerate}
 \item For any specified nonempty finite set $U$, there is a graph on
 $U$ in which every two disjoint sets $X,Y\subseteq U$ with
 $|X|,|Y|\geq\lceil\zeta|U|\rceil$ have an edge between them.
 \item For any specified disjoint nonempty finite sets $A,B$ with
 $\max\{|A|,|B|\}/\min\{|A|,|B|\}\leq\kappa$, there is a bipartite
 graph on $A\mathbin{\dot\cup}B$ with parts $A,B$ in which every
 $X\subseteq A$, $Y\subseteq B$ with
 $|X|,|Y|\geq\lceil\zeta(|A|+|B|)\rceil$ have an edge between them.
\end{enumerate}
\end{restatedlemma}

\begin{proof}[Proof of Lemma~\ref{lem:factor-expanders}]
Each auxiliary vertex has $d=8^h$ outgoing and incoming steps,
because its adjacency matrix is symmetric with row sum $d$.
For any output vertex $v$, let $\rho$ be its map ($\pi$, $\pi_A$,
or $\pi_B$), and let $n_\rho$ be the size of that map's prescribed
set. Count outgoing steps from its preimage, except on side $B$,
where we count incoming steps. Each output neighbor uses at least
one such step, so the same bound covers all three cases:
\[
 \deg(v)\le8^h|\rho^{-1}(\{v\})|
 \overset{\eqref{eq:factor-expander-preimages}}{\leq}
 8^h\lceil M/n_\rho\rceil.
\]

In the one-set construction,
\begin{equation}\label{eq:factor-expander-order}
 M=\lceil\sqrt n\rceil^2\le(\sqrt n+1)^2\le4n.
\end{equation}
Hence $\lceil M/n\rceil\le4$. For the bipartite construction, the
same estimate with $n=\max\{n_A,n_B\}$ gives
$M\overset{\eqref{eq:factor-expander-order}}{\leq}4\max\{n_A,n_B\}$.
The side-size ratio
is at most $\kappa$, so
$M/n_A,M/n_B\le4\kappa$. Both constructions therefore have degree
at most $\lceil4\kappa\rceil8^h=d_{\zeta,\kappa}$.

We next verify the no-hole conditions at the thresholds in the
lemma. In the one-set case, put $t=\lceil\zeta n\rceil$. Then
$\lambda_0^hM<(\zeta/4)M
\overset{\eqref{eq:factor-expander-order}}{\leq}\zeta n\le t$.
In the bipartite case, put $t=\lceil\zeta(n_A+n_B)\rceil$. Then
$\lambda_0^hM<(\zeta/4)M
\overset{\eqref{eq:factor-expander-order}}{\leq}\zeta\max\{n_A,n_B\}\le t$.
We apply the mixing bound to preimages of $t$-sets. If no such admissible sets exist, the corresponding
no-hole condition is vacuous. Otherwise, any larger admissible sets
contain admissible $t$-element subsets.
For the one-set construction, take disjoint $t$-sets $X,Y\subseteq U$
and put $S=\pi^{-1}(X)$, $T=\pi^{-1}(Y)$.
For the bipartite construction, take $t$-sets $X\subseteq A$,
$Y\subseteq B$ and put $S=\pi_A^{-1}(X)$, $T=\pi_B^{-1}(Y)$.
In both cases surjectivity gives $|S|,|T|\ge t>\lambda_0^hM$,
so \eqref{eq:factor-expander-mixing} supplies a step from $S$ to $T$.
Its endpoint images give an edge between $X$ and $Y$. In the one-set
case the preimages are disjoint and their images are distinct.
In the bipartite case the preimages may overlap, and even a loop
step gives a valid edge because its images lie in the disjoint
sets $A,B$.

It remains to prove connectedness. If $\mathcal E$ had a nonempty
proper component $C$, then
$x=\mathbf1_C-(|C|/M)\mathbf1$ would be a nonzero zero-sum vector
with $A_{\mathcal E}x=dx$, contradicting its nonconstant eigenvalue
bound $d\lambda_0^h<d$. Thus $\mathcal E$ is connected.
It also has a loop at $(0,0)$: repeat $h$ times a loop of the
Margulis--Gabber--Galil graph at that vertex. This explicitly
supplies the odd closed walk needed below.

In the one-set construction, applying $\pi$ to an auxiliary walk
gives a sequence of equal or adjacent output vertices. Deleting
consecutive repetitions gives an output walk. Connectedness of
$\mathcal E$ and surjectivity of $\pi$ thus prove connectedness of
the output.

For the bipartite construction we need auxiliary walks of both
parities. From any $x\in V(\mathcal E)$, walk to $(0,0)$, traverse its loop, and
return along the reversed walk. Prepending this odd closed walk to a walk from
$x$ to $y\in V(\mathcal E)$ changes its parity. Map a walk
$x=x_0,x_1,\ldots,x_\ell=y$ to the alternating sequence
$\pi_A(x_0),\pi_B(x_1),\pi_A(x_2),\ldots$ in the output.
Each consecutive pair is an edge by construction, using the
reverse auxiliary step when the sequence goes from $B$ to $A$.
An even-length walk reaches $\pi_A(y)$, and an odd-length walk
reaches $\pi_B(y)$. Thus $\pi_A(x)$ is connected to both images of
every $y$, and surjectivity gives connectedness.

Both outputs are simple by their edge definitions.
The auxiliary order $M$ is at most four times the largest prescribed
set size. Theorem~\ref{thm:auxiliary-expander-family} constructs its
walk-power graph in polynomial time for fixed $h$. Map the endpoints
as specified above and sort the resulting pairs to merge repeated
edges, discarding equal images in the one-set construction. Since
$h$ depends only on $\zeta$, these operations take polynomial time
for the fixed parameters.
\end{proof}
